% ============================================================================
% Fichier  : partie4/P04_C12_securite_donnees.tex
% Titre    : Auditer la Sécurité des Données
% Auteur   : MINKA MI NGUIDJOÏ Thierry Emmanuel
% Version  : 1.0 -- novembre 2026
% Corrigés : corriges/Corr_C12.tex
% ============================================================================
% Sources : loi n° 2024/017, art. 7, 15, 22(2), 24, 25(2-3), 27, 28, 57
% (texte officiel) ; ISO/IEC 27002:2022 et ISO/IEC 27701 (mesures,
% paraphrasées).
% Faits nouveaux (bible) : consultations du dossier d'une patiente les 11 et
% 12 mai 2026 au CENTRE MÉDICAL MBOLO ; 14 consultations, 6 comptes, 4 sans
% lien avec la prise en charge.
% ============================================================================

\chapter{Auditer la Sécurité des Données}
\label{chap:securite}
\chaptermark{La sécurité des données}

\epigraph{%
    « La menace la plus fréquente sur un dossier médical\\
    n'est pas le pirate. C'est le collègue curieux. »%
}{Principe du cours}

\niveauChapitre{Intermédiaire}{%
    Chapitres~\ref{chap:audit} et~\ref{chap:risque}~; bases de la sécurité
    des systèmes
}

\begin{boxobjectifs}
\begin{itemize}
    \item Énoncer les obligations de sécurité des articles~22, 24, 25
          et~27 de la loi.
    \item Traduire les dix exigences de l'article~27 en contrôles
          vérifiables.
    \item Exploiter des journaux d'accès pour établir des consultations
          sans lien avec la finalité.
    \item Distinguer la sécurité des systèmes de la protection des données
          contre les usages internes illégitimes.
    \item Auditer le rapport annuel de l'article~27, alinéa~2.
\end{itemize}
\end{boxobjectifs}

\bigskip

% ============================================================================
\section*{Avant de commencer : quatorze consultations en quarante-huit heures}
% ============================================================================

Bertoua, mercredi 13~mai~2026. Une chanteuse connue de la région, admise
deux jours plus tôt au CENTRE MÉDICAL MBOLO pour une intervention, apprend
par un message vocal qui circule en ville le nom de son affection. Elle
appelle la directrice. Le Dr~ABENA demande à M.~YOUSSOUFA d'extraire du
dossier patient informatisé la liste des consultations de ce dossier.

L'extraction tient sur une page~: quatorze consultations en quarante-huit
heures, par six comptes. Deux sont ceux du chirurgien et de l'infirmière du
service. Les quatre autres sont ceux d'une secrétaire de la maternité, d'un
agent de la facturation, d'un brancardier dont le compte est partagé par
l'équipe de nuit, et d'un stagiaire dont le stage s'est terminé en mars.

\begin{boxalert}
\textbf{Violation typique.} Auditer la sécurité des données comme la sécurité
d'un réseau~: pare-feu, antivirus, mots de passe, sans regarder qui, à
l'intérieur, consulte quoi.

\smallskip
\textbf{Règle~:} la loi exige que les personnes autorisées n'accèdent
« qu'aux données relevant de leur autorisation » et que l'on puisse
vérifier « a posteriori » qui a consulté quoi (article~27). L'audit teste
les droits d'accès et exploite les journaux.
\end{boxalert}

% ============================================================================
\section{Les Obligations de Sécurité}
\label{sec:c07-obligations}
% ============================================================================

Les obligations de sécurité sont dispersées dans la loi. Les articles~7
et~15 posent la confidentialité~; l'article~24 réserve le traitement aux
personnes « qui agissent sous l'autorité du responsable du traitement et
seulement sur ses instructions », et renvoie au référentiel de mesures de
l'Autorité~; l'article~25, alinéa~2, impose d'empêcher que les données
soient « déformées, endommagées, indisponibles » ou que des tiers non
autorisés y aient accès~; l'article~22, alinéa~2, cite le chiffrement et
« l'instauration des autorisations restrictives d'accès aux données »~; l'article~27 en
donne la liste la plus complète.

\begin{boxjuris}
\textbf{\loicam{n°~2024/017}, article~27, alinéa~1 (extraits)}~\cite{loi2024017}~:
\textit{« Le responsable du traitement et le sous-traitant des données à
caractère personnel sont tenus~: de garantir la disponibilité des données à
caractère personnel notamment grâce à la mise en œuvre des mesures visant à
garantir la résilience des systèmes d'information contenant les données à
caractère personnel~; de garantir que lors de l'utilisation d'un système de
traitement automatisé de données, les personnes autorisées ne puissent
accéder qu'aux données relevant de leur autorisation~; [\ldots] de
garantir que puisse être vérifiée et constatée a posteriori l'identité des
personnes ayant accès au système d'information contenant des données à
caractère personnel, la nature des données qui ont été introduites,
modifiées, altérées, copiées, effacées ou lues dans le système, ainsi que
la période de manipulation de ces données~; [\ldots] de sauvegarder les
données par la constitution de copies de sécurité protégées~; [\ldots] »}

\smallskip
\textbf{Portée.} L'article énumère dix obligations. La plus exigeante pour
l'audit est la traçabilité~: la loi vise les modifications, et aussi les
données « lues ». Un système qui ne journalise pas les
consultations ne permet pas de satisfaire l'article.
\end{boxjuris}

L'alinéa~2 du même article ajoute une obligation de compte rendu~: le
responsable « adresse à l'Autorité de protection des données à caractère
personnel un rapport annuel sur l'état de mise en œuvre des mesures de
sécurité contenues dans le référentiel des mesures techniques et
organisationnelles ». Tant que le référentiel n'est pas publié, le contenu
du rapport est un \NEK{} de critère~; l'existence d'un état annuel des
mesures, elle, peut être vérifiée.

% ============================================================================
\section{Dix Obligations, Dix Contrôles}
\label{sec:c07-controles}
% ============================================================================

\begin{tableaucours}{L{5.0cm}L{3.0cm}L{4.0cm}}{%
    \tableauHeader{Obligation (article~27)} &
    \tableauHeader{Mesure ISO/IEC 27002} &
    \tableauHeader{Test d'audit}
}
Disponibilité et résilience & 5.30, 8.14 & Plan de reprise exercé, durée
d'indisponibilité mesurée \\
Accès limité à l'autorisation & 5.15, 5.18, 8.3 & Comparer droits et
fonctions sur échantillon \\
Identité des destinataires vérifiée & 5.14 & Procédure de transmission,
échantillon d'envois \\
Traçabilité a posteriori & 8.15 & Journal des consultations, test de
reconstitution \\
Accès physique aux installations & 7.1, 7.2 & Visite, registre des accès \\
Supports protégés & 7.10 & Inventaire, chiffrement des supports amovibles \\
Introduction et modification non autorisées & 8.2, 8.5 & Comptes
privilégiés, validation des saisies \\
Transmission protégée & 5.14, 8.24 & Chiffrement des flux, configuration \\
Sauvegardes protégées & 8.13 & Test de restauration, sauvegarde hors
ligne \\
Autres mesures réglementaires & À préciser & \NEK{} de critère \\
\end{tableaucours}
\vspace{-0.8em}
\begin{center}{\small\itshape Tableau~12.1. L'article~27 et les mesures de l'ISO/IEC~27002:2022 (correspondances indicatives)}\end{center}

Les numéros de mesures renvoient à l'ISO/IEC~27002:2022~\cite{iso27002_2022}
et sont donnés à titre indicatif~; la correspondance n'est pas fixée par la
loi. L'ISO/IEC~27701~\cite{iso27701} complète ces mesures par des exigences
propres aux données personnelles (Chapitre~\ref{chap:normes}).

\subsection{Les droits d'accès}

L'audit des droits d'accès compare, pour un échantillon de comptes, les
droits accordés aux besoins de la fonction. Trois anomalies reviennent
partout~: les comptes de personnes parties (stagiaires, intérimaires,
salariés licenciés), les comptes partagés par une équipe, qui rendent la
traçabilité impossible, et les droits trop larges accordés « pour
dépanner » et jamais retirés.

\subsection{Les journaux de consultation}

Exploiter un journal de consultation, c'est répondre à une question précise~:
« Qui a consulté ce dossier, et chaque consultation avait-elle un lien avec
la finalité~? » On rapproche chaque ligne du journal d'une pièce qui la
justifie~: le planning du service, l'affectation du patient, l'acte facturé.
Une consultation sans pièce justificative n'est pas encore une faute~; c'est
une consultation à expliquer.

\begin{boxinfo}
\textbf{Un journal se lit avec ses limites.} Un compte partagé ne désigne pas
une personne~; une consultation par un compte ne prouve pas qui était au
clavier~; un journal qui ne garde que trente jours ne permet pas de
remonter plus loin. L'auditeur écrit ce que le journal établit (« le compte
X a ouvert le dossier à telle heure ») et ce qu'il n'établit pas (« la
personne derrière le compte »).
\end{boxinfo}

% ============================================================================
\section{Le Dossier de la Patiente}
\label{sec:c07-mbolo}
% ============================================================================

\begin{boxscenario}{Mission MBOLO. Les quatorze consultations, 11 et 12 mai 2026}
\begin{tableaucours}{L{3.0cm}cL{3.2cm}L{3.6cm}}{%
    \tableauHeader{Compte} &
    \tableauHeader{Consult.} &
    \tableauHeader{Fonction} &
    \tableauHeader{Pièce justificative}
}
Chirurgien & 4 & Bloc opératoire & Programme opératoire \\
Infirmière du service & 4 & Chirurgie & Planning, soins notés \\
Secrétaire maternité & 2 & Maternité & Aucune \\
Agent de facturation & 1 & Facturation & Facture émise le 12~mai \\
Compte « nuit\_brancard » & 2 & Partagé, 5 brancardiers & Aucune \\
Stagiaire & 1 & Stage terminé en mars & Aucune \\
\end{tableaucours}

\medskip
\textbf{Constatations.} Quatorze consultations. Neuf sont justifiées par une
pièce (chirurgien, infirmière, facturation). Cinq ne le sont pas~: deux par
la secrétaire de la maternité, deux par un compte partagé entre cinq
personnes, une par le compte d'un stagiaire parti depuis deux mois.

\textbf{Analyse} (confiance élevée sur les comptes, faible sur les
personnes). Cinq consultations sont sans lien établi avec la prise en
charge. Le journal ne permet pas de savoir lequel des cinq brancardiers a
utilisé le compte partagé, ni qui a utilisé le compte du stagiaire. La
proposition « la fuite vient d'une consultation du dossier » reste \NEK{}~:
le diagnostic a pu circuler oralement.
\end{boxscenario}

Le Dr~ABENA a désactivé le compte du stagiaire et le compte partagé le jour
même, et ouvert une procédure disciplinaire pour la secrétaire, qui a
reconnu la consultation. Elle a aussi posé à KORO AUDIT, mandaté en
octobre, une question plus large~: combien d'autres dossiers ont été
consultés de la même façon sans que personne ne s'en plaigne~?

\begin{boxlecture}{Le domaine D3 à MBOLO, octobre 2026}
\begin{tableaucours}{L{7.8cm}cL{3.2cm}}{%
    \tableauHeader{Exigence} &
    \tableauHeader{Cote} &
    \tableauHeader{Preuve}
}
Les consultations sont journalisées & \OUI{} & Journal MEDIKA \\
Les journaux sont conservés au moins un an & \NON{} & 90 jours \\
Chaque compte désigne une seule personne & \NON{} & 3 comptes partagés \\
Les comptes des personnes parties sont désactivés & \NON{} & 7 comptes
actifs sur 31 départs \\
Les droits correspondent aux fonctions & \NON{} & Facturation voit les
diagnostics \\
Les journaux sont revus & \NON{} & Aucune revue avant mai \\
Les sauvegardes sont restaurables & \OUI{} & Test de restauration du
22~octobre \\
Les flux de téléconsultation sont chiffrés de bout en bout & \NON{} &
Option non activée \\
Le prestataire applique les mesures convenues & \NEK{} & MEDISOFT non
audité \\
Le rapport annuel de l'article~27(2) est conforme & \NEK{} & Référentiel non
publié \\
\end{tableaucours}

\medskip
\textbf{Décompte.} Deux \OUI{}, six \NON{}, deux \NEK{} sur dix.
\hfill\lecture{0{,}20}{0{,}20}{1{,}37}

\medskip
\textbf{Lecture.} Domaine opaque au seuil ($U = 0{,}20$), mais six écarts
établis qui touchent tous la traçabilité et les droits d'accès. La
journalisation existe~; c'est son usage qui manque.
\end{boxlecture}

% ============================================================================
\section{La Sécurité à NYANGA, après l'Incident}
\label{sec:c07-nyanga}
% ============================================================================

À NYANGA CRÉDIT, la sécurité s'est renforcée depuis mars~2025~: la
passerelle a été remplacée, l'inventaire des équipements est tenu, les
journaux sont centralisés. L'audit le constate et cote \OUI{} les
exigences correspondantes. Il relève aussi ce que l'incident n'a pas
changé~: les agents des agences voient les comptes de tous les clients du
réseau, et non de leur seule agence~; les relevés de compte partent encore
par messagerie personnelle quand un client le demande. La sécurité
technique a progressé~; la limitation des accès à la finalité, non.

% ============================================================================
\section{Ce que la Sécurité demande à chaque Métier}
\label{sec:c07-metiers}
% ============================================================================

\begin{boxmetier}{RSSI}
Supprimer les comptes partagés, désactiver les comptes à la date de départ,
conserver les journaux de consultation assez longtemps pour enquêter, et
les revoir régulièrement sur les dossiers sensibles.
\end{boxmetier}

\begin{boxmetier}{Auditeur de la protection des données}
Tester les droits d'accès contre les fonctions et rapprocher les journaux de
pièces justificatives. Écrire ce que le journal établit sur un compte, et ce
qu'il n'établit pas sur une personne.
\end{boxmetier}

\begin{boxmetier}{Responsable de traitement}
Tirer les conséquences disciplinaires d'une consultation injustifiée, et
faire savoir au personnel que les consultations sont tracées et revues~:
la dissuasion fait partie de la mesure.
\end{boxmetier}

\begin{boxresume}
\begin{itemize}
    \item Les obligations de sécurité figurent aux articles~7, 15, 22, 24,
          25 et~27~; l'article~27 en donne dix, dont la traçabilité des
          données lues.
    \item Le responsable adresse un rapport annuel sur les mesures à
          l'Autorité (article~27, alinéa~2).
    \item Les menaces internes se contrôlent par les droits d'accès et
          l'exploitation des journaux de consultation.
    \item Un journal établit ce qu'un compte a fait, pas ce qu'une personne a
          fait~; les comptes partagés détruisent la traçabilité.
    \item À MBOLO, cinq consultations d'un dossier sensible sont sans lien
          établi avec la prise en charge~; six écarts sont établis dans le
          domaine de la sécurité.
\end{itemize}
\end{boxresume}

\begin{boxinfo}
\textbf{Points de contrôle.} Vous avez compris ce chapitre si vous savez~:
\begin{itemize}
    \item traduire une obligation de l'article~27 en test d'audit~;
    \item rapprocher un journal de consultation de pièces justificatives~;
    \item dire pourquoi un compte partagé est un écart à l'article~27.
\end{itemize}
\end{boxinfo}

\begin{boxexo}[title={\ding{45}\quad Exercice~12.1. Quel contrôle \niveaubase}]
Proposez un test d'audit pour chaque obligation~: (a) les sauvegardes sont
protégées~; (b) les supports amovibles ne sont pas lisibles par un tiers~;
(c) l'identité des destinataires est vérifiée avant transmission~;
(d) les données sont disponibles en cas de panne.
\end{boxexo}

\begin{boxexo}[title={\ding{45}\quad Exercice~12.2. Lire un journal \niveaumoyen}]
Le compte de l'agent de facturation a consulté le dossier une fois, le
12~mai à 16~h~04~; la facture a été émise à 16~h~11. Le compte a aussi accès
à l'onglet « diagnostic ». Quelles cotes donnez-vous aux propositions
« cette consultation était justifiée » et « les droits de ce compte sont
limités à sa fonction »~?
\end{boxexo}

\begin{boxexo}[title={\ding{45}\quad Exercice~12.3. Combien d'autres dossiers \niveaumoyen}]
Proposez une méthode pour répondre à la question du Dr~ABENA~: combien
d'autres dossiers ont été consultés sans lien avec la prise en charge~?
Précisez les données à extraire, les règles de rapprochement, l'échantillon,
et la manière de protéger les données des patients pendant l'analyse.
\end{boxexo}

\begin{boxexo}[title={\ding{45}\quad Exercice~12.4. Sécurisé mais non conforme \niveauexpert}]
Le DSI de NYANGA affirme~: « Depuis l'incident, notre sécurité est au
niveau des meilleures banques. » Montrez, à partir des constats de la
section~\ref{sec:c07-nyanga}, en quoi cette affirmation peut être vraie et
l'établissement non conforme à l'article~27. Rédigez la recommandation.
\end{boxexo}

\bigskip

\begin{boxinfo}
\textbf{Pour aller plus loin.}
\begin{itemize}
    \item L'ISO/IEC~27002:2022~\cite{iso27002_2022}, mesures de contrôle
          d'accès et de journalisation.
    \item L'ISO/IEC~27701~\cite{iso27701}, pour les mesures propres aux
          données personnelles.
\end{itemize}
\end{boxinfo}

\begin{boxinfo}
\textbf{Transition.} La Partie~II a montré comment auditer. Il faut
maintenant lire de près le cadre qui fournit les critères~: la loi
n°~2024/017 elle-même, son Autorité, ses sanctions, et le cadre africain
dans lequel elle s'inscrit. C'est l'objet du Chapitre~\ref{chap:loi}.
\end{boxinfo}
